% Polished version of intro-xiao.tex (proofread by Claude). The original file is unchanged.
\section{Introduction}
Spatial reasoning and geometric understanding are fundamental to human cognition and to artificial intelligence, supporting applications from robotics and navigation to computer-aided design. As large language models (LLMs) have matured, a growing body of work asks whether they can reason about shape, position, and transformation. Recent benchmarks probe this ability through paper folding and unfolding~\cite{zhou2026paperfoldingpuzzle} and through geometry expressed as procedural drawing code~\cite{luo2026geogrambench}, and both report pronounced deficiencies in current models. These studies largely assess spatial reasoning by asking a model to select or compute an answer. A complementary and less examined question is whether a language model can \emph{construct}: given only words, can it produce a three-dimensional object whose parts are correctly sized, correctly placed, and physically joined?

LEGO brick assembly makes this question concrete. A concise request such as \emph{a simple chair} identifies an object category but leaves its dimensions, proportions, and part layout largely unspecified. Answering it requires three coupled forms of geometric spatial reasoning. The model must complete the missing geometry from prior knowledge of what the object looks like. It must discretize that geometry into pieces of fixed dimensions at integer positions on a grid. And it must respect the difference between spatial adjacency and stud connectivity: pieces that touch side by side do not hold together unless they interlock vertically through studs. We call this setting the \texttt{Text-to-Brick} task: translating a natural-language description, at any level of detail, into a set of LEGO brick placements that reproduces the described shape as a single connected assembly. Because every output can be checked for coverage, bounds, and connectivity, the task offers a verifiable test of whether linguistic knowledge has become spatial structure.

Prior work has shown that language models can be trained for this task. BrickGPT~\cite{brickgpt} learns autoregressive brick-placement sequences and applies validity checks with physics-aware rollback during inference. BrickNet~\cite{bricknet} adopts a graph-based representation that supports a broader range of parts and connection types. Two questions remain open. \textbf{(Q1)} How well do generators translate \emph{limited} linguistic information into object geometry? Pooled scores hide this: only a level-wise evaluation (Table~\ref{tab:level-results}) shows, for example, that BrickNet trails BrickForma on mean VQAScore, PE, and SigLIP2 even at the two shortest description levels, L1 and L2. \textbf{(Q2)} Can LLM-assisted programs encode the geometric details omitted from concise prompts in a form suitable for training a compact model?

We present BrickForma, which uses executable programs to generate training examples for a compact text-to-brick model. With the assistance of GPT-6.1 Sol, we develop parametric programs that vary object proportions, part dimensions, openings, and support arrangements to synthesize voxel shapes in batches. Each program specifies the structure of an object family and exposes its geometric degrees of freedom as parameters. The source collection contains 80,000 models across 21 categories on a $48^3$ grid. Each model is paired with five levels of description that progressively add structural, geometric, proportional, and appearance information. The resulting approximately 400,000 caption records link concise requests and detailed specifications to shared geometry, providing supervision for prompts with different amounts of geometric detail.

Turning voxel shapes into brick targets is where spatial adjacency and stud connectivity must be separated: pieces can cover a connected voxel volume and touch side by side without forming any stud connection. Our conversion pipeline combines same-color brick placement, connection analysis, and local refinement to construct geometrically checked target layouts, and it records every change to the source voxels in the derived models. The integrated dataset release contains 79,301 models that pass its dataset-level geometric checks. We fine-tune Qwen3-1.7B~\cite{qwen3} with low-rank adaptation~\cite{lora} on caption--placement pairs to predict brick dimensions and positions.

Geometric refinement complements the small model's predictions with explicit structural enforcement. From the supported, in-bounds portions of the predicted placements, the pipeline forms a voxel union and uses mixed-integer linear programming to find an exact-cover layout that forms a single stud-connected component. Independent coverage and connectivity checks then verify the refined layout. This design gives each half of the task to the component suited to it: the language model decides what shape a description denotes, and the solver decides how that shape holds together.

We evaluate BrickForma against BrickGPT and BrickNet on 525 shared prompts spanning five description levels. BrickForma's final outputs after refinement and recovery achieve the highest mean VQAScore, PE, and SigLIP2 at every level, including the shortest descriptions. The mean invalid-placement ratio is 0.00\% for BrickForma and BrickGPT, and 18.42\% for BrickNet. Reporting all five levels links the multi-level data design to an explicit assessment of how generation changes with the amount of descriptive detail.

In summary, our main contributions are as follows:

\textbf{LLM-assisted training data for a compact model.}
We use executable parametric programs to generate voxel shapes and their descriptions, convert the shapes into geometrically checked brick layouts, and fine-tune a 1.7B-parameter model on the resulting caption--placement pairs. On the 525 shared prompts, BrickForma's final outputs reach a mean VQAScore of 0.0872, compared with 0.0666 for BrickGPT, an absolute difference of 0.0206 (30.9\% relative).

\textbf{A multi-level dataset for evaluating alignment across levels of descriptive detail.}
Our pipeline generates 80,000 source voxel models with five description levels each; in the integrated release, 79,301 of them yield brick layouts that pass the geometric checks. Level-wise evaluation against BrickGPT and BrickNet, covering text--image alignment and invalid placements from concise prompts to detailed specifications, shows that BrickForma achieves the highest mean alignment scores at all five levels.
